% Exact LaTeX recreation of the supplied Phase-1.pdf
% Compile with pdfLaTeX. Place member pictures in the same folder if required.

\documentclass[11pt,a4paper]{article}

\usepackage[a4paper,left=2.0cm,right=2.0cm,top=1.55cm,bottom=1.55cm]{geometry}
\usepackage[T1]{fontenc}
\usepackage{lmodern}
\usepackage{graphicx}
\usepackage{xcolor}
\usepackage{array}
\usepackage{tabularx}
\usepackage{ragged2e}
\usepackage{enumitem}
\usepackage{fancyhdr}
\usepackage{tikz}
\usepackage{setspace}
\usepackage{graphicx}

% -------------------------------------------------
% COLOURS
% -------------------------------------------------
\definecolor{TitleBlue}{RGB}{61,103,154}

% -------------------------------------------------
% GENERAL SETTINGS
% -------------------------------------------------
\setlength{\parindent}{0pt}
\setlength{\parskip}{0pt}
\renewcommand{\arraystretch}{1.0}

% The supplied document uses a clean sans-serif heading style
% and italic serif-style explanatory text.
\newcommand{\blueheading}[1]{%
    {\sffamily\color{TitleBlue}\fontsize{15}{18}\selectfont #1}%
}

\newcommand{\bluebigheading}[1]{%
    {\sffamily\bfseries\color{TitleBlue}\fontsize{16}{19}\selectfont #1}%
}

\newcommand{\instruction}[1]{%
    {\itshape\fontsize{10.5}{12.5}\selectfont #1}%
}

\newcommand{\answerbox}[2][3.0cm]{%
    \vspace{2pt}
    \noindent
    \fbox{%
        \begin{minipage}[t][#1][t]{\dimexpr\textwidth-2\fboxsep-2\fboxrule\relax}
            \vspace{1mm}
            \itshape #2
        \end{minipage}
    }
}

% -------------------------------------------------
% HEADER / FOOTER
% -------------------------------------------------
\pagestyle{fancy}
\fancyhf{}
\renewcommand{\headrulewidth}{0pt}
\renewcommand{\footrulewidth}{0pt}

\fancyhead[L]{%
    \sffamily\bfseries\itshape\fontsize{10.5}{12}\selectfont
   DSCPP-III-2026-T382
}
\fancyhead[R]{%
    \sffamily\bfseries\fontsize{10.5}{12}\selectfont
    PROJECT PROPOSAL
}
\fancyfoot[R]{%
    \itshape\fontsize{10}{12}\selectfont
    Page \thepage
}
\setlength{\headheight}{14pt}
\setlength{\headsep}{23pt}
\setlength{\footskip}{24pt}

% -------------------------------------------------
% PAGE 1
% -------------------------------------------------
\begin{document}

\vspace*{4mm}

\bluebigheading{PROJECT AND TEAM INFORMATION}

\vspace{13mm}

\blueheading{Project Title}

\vspace{1mm}

\instruction{}

\answerbox[1.0cm]{\textbf{\textit{HABITQUEST}}: Gamified Habit Tracker }

\vspace{13mm}

\blueheading{Student / Team Information}

\vspace{2mm}

% Team information table
% Compact 4-member layout so all members remain on Page 1.
\noindent
\begin{tabularx}{\textwidth}{|>{\RaggedRight\arraybackslash}p{0.69\textwidth}|>{\RaggedRight\arraybackslash}X|}
\hline
\begin{minipage}[t]{\linewidth}
\vspace{1.5mm}
\textit{Team Name:}\\[-1mm]
\textit{Team ID:}
\vspace{1.5mm}
\end{minipage}
&
\begin{minipage}[t]{\linewidth}
\vspace{1.5mm}
\textbf{\textit{Alpha\\DSCPP-III-2026-T382}}
\vspace{1.5mm}
\end{minipage}
\\
\hline
\begin{minipage}[t][3.25cm][t]{\linewidth}
\vspace{1.5mm}
\textbf{Team member 1 (Team Lead)}\\[-1mm]
\vspace{2.0cm}
\end{minipage}
&
\begin{minipage}[t][3.25cm][t]{\linewidth}
\vspace{1.5mm}
\textit{Semwal, Ayush -- 25101210653}\\
\resizebox{\linewidth}{!}{\textit{AYUSHSEMWALSEMWAL.25101210653@gehu.ac.in}}

\vspace{2mm}
\includegraphics[width=3.0cm,height=1.5cm,keepaspectratio]{ayushphoto.jpg}
\end{minipage}
\\
\hline
\begin{minipage}[t][3.25cm][t]{\linewidth}
\vspace{1.5mm}
\textbf{Team member 2}\\[-1mm]

\vspace{2.0cm}
\end{minipage}
&
\begin{minipage}[t][3.25cm][t]{\linewidth}
\vspace{1.5mm}
\textit{Pratap Singh , Tarun -- 25101210657}\\
\resizebox{\linewidth}{!}{\textit{TARUNPRATAPSINGH.25101210657@gehu.ac.in}}

\vspace{2mm}
\includegraphics[width=3.0cm,height=1.5cm,keepaspectratio]{tarunphoto.jpg}
\end{minipage}
\\
\hline
\begin{minipage}[t][3.25cm][t]{\linewidth}
\vspace{1.5mm}
\textbf{Team member 3}\\[-1mm]
\vspace{2.0cm}
\end{minipage}
&
\begin{minipage}[t][3.25cm][t]{\linewidth}
\vspace{1.5mm}
\textit{ Rawat, Suryansh -- 25101210654}\\
\resizebox{\linewidth}{!}{\textit{SURYANSHRAWAT.25101210654@gehu.ac.in}}

\vspace{2mm}
\includegraphics[width=3.0cm,height=1.5cm,keepaspectratio]{suryanshphoto.png}
\end{minipage}
\\
\hline
\begin{minipage}[t][3.25cm][t]{\linewidth}
\vspace{1.5mm}
\textbf{Team member 4}\\[-1mm]
\vspace{2.0cm}
\end{minipage}
&
\begin{minipage}[t][3.25cm][t]{\linewidth}
\vspace{1.5mm}
\textit{
Bisht, Priyanshu -- 25101281855}\\
\resizebox{\linewidth}{!}{\textit{PRIYANSHUBISHT.25101281855@gehu.ac.in}}

\vspace{2mm}
 \includegraphics[width=3.0cm,height=1.5cm,keepaspectratio]{priyanshu.jpg}
\end{minipage}
\\
\hline
\end{tabularx}

\newpage

% -------------------------------------------------
% PAGE 2
% -------------------------------------------------

\bluebigheading{PROPOSAL DESCRIPTION (10 pts)}

\vspace{12mm}

\blueheading{Motivation (1 pt)}

\vspace{1mm}


\answerbox[3.3cm]{The main problem that we are aiming to solve with this project is the difficuilty of maintaining habits with consistency over time.
The traditional habit trackers mostly work only as simple checklist and does not provide motivation, but our gamified habit tracker solves this problem through a reward based system where users earns XP for completing daily habits, receive streak bonus and can unlock new levels and achievement badges, making habit tracking more engaging and fun by providing a continuous feedback and a sense of progress . our system will encourage users to stay consistent and improve their habits, making their lives more productive and healthy.
 }

\vspace{12mm}

\blueheading{State of the Art / Current solution (1 pt)}

\vspace{1mm}


\answerbox[2.75cm]{Most of the gamified habit trackers uses heavy resources GUI or mobile applications but this project replaces those heavy graphical interfaces with a fast, terminal based CLI application. Instead of heavy libraries it uses standard input /output streams and ANSI escape  sequences to provide user coloured visual feedback , streak counter and dynamic health bars within the command line hence making this project light weight /fast but at the same time more interactive and appealing to the users.

}

\vspace{12mm}

\blueheading{Project Goals and Milestones (2 pts)}

\vspace{1mm}


\answerbox[8.15cm]{The main goal of our project is to help people build discipline, develop positive habits, increase productivity and become the best version of themselves through a simple, light and engaging gamified habit tracker . We aim to motivate users through XP, rewards, streaks, new levels and acheivements. Initial milestones
Include designing the system , applying habit tracking and reward logic, developing the CLI interface and testing core features. Later milestones include adding analytics, improving gamification, optimizing the system and performance and managing user testing. In the long term we aim to scale the project into a widely accessible platform that can and will help people build disciplined and an active lifestyle. The primary goal is to build a fully functioning C/C++ application that keeps you on track through interactive gamification.\\
\\The main steps/approach to design the project is as follows:
\\1: Define core requirements and design OOPs class hierarchies .
\\2: Implement C/C++ data structures including dynamic memory allocation and BSTs(binary search trees) for leaderboards.
\\3: Develop the gamification engine (streak logic, penalties,XP calculation).
\\4 : Build an interactive CLI loop and apply ANSI color formatting.
\\5: Integrate local flat file serialization and conduct memory leak testing.      
}

\newpage 
\vspace{12mm}

\blueheading{Project Approach (3 pts)}

\vspace{1mm}


\answerbox[6.1cm]{We have planned to make this Habit Tracker a simple system that will help users stay consistent with their daily habits and routines . The targeted user will be able to create habits , mark them as completed , list them and keep track of their progress through a command line interface. 
\\To make the tracker we are adding a gamification system where the user will earn XP(experience points) for completing their daily habits, maintain streaks, level up and unlock new achievement badges.maintaining Streaks will provide extra XP to the users so that they extra motivation to stay conssistent.The application will also display some pop up coloured alerts and simple statistics which will help users understand their progress in a better and easy way. A special feature of heart or health bar will add an extra layer of challenge for the users as it will get depleted if the habits are missed.
\\ All important information like habits, XP,Streaks, Levels and achievements will be stored localy by using .dat files, so that the user’s progress remains available even after closing the program.  We will use C++, OOPs concepts and data structures to keep the  project organised and make it easier for future improvements and expanssion. 

}


% -------------------------------------------------
% PAGE 3
% -------------------------------------------------
\vspace{12mm} 

\blueheading{System Architecture (High Level Diagram)(2 pts)}

\vspace{1mm}


\answerbox[4.0cm]{{\textit{\textbf{Flow chart of the System Architecture :}}}
\usetikzlibrary{positioning,arrows.meta}

\begin{center}
\begin{tikzpicture}[
    node distance=4mm and 4mm,
    box/.style={
        draw,
        rounded corners=2pt,
        align=center,
        text width=2.45cm,
        minimum height=0.85cm,
        inner sep=3pt,
        font=\scriptsize\sffamily
    },
    arrow/.style={
        -{Latex[length=1.5mm]},
        semithick
    }
]

\node[box] (input) {
    \textbf{Input Layer}\\
    CLI Parser\\
    Habit Logs\\
    Menu Navigation
};

\node[box, right=of input] (controller) {
    \textbf{Activity Controller}\\
    OOP Classes\\
    Timestamp Validation\\
    Streak Reset
};

\node[box, right=of controller] (game) {
    \textbf{Gamification Engine}\\
    XP Calculation\\
    Streak Multipliers\\
    Heart/Health Deduction
};

\node[box, right=of game] (storage) {
    \textbf{Storage Module}\\
    Serialization\\
    .dat Flat-Files\\
    Data Persistence
};

\node[box, right=of storage] (output) {
    \textbf{Output Layer}\\
    Terminal Renderer\\
    ANSI Alerts\\
    ASCII Statistics
};

\draw[arrow] (input) -- (controller);
\draw[arrow] (controller) -- (game);
\draw[arrow] (game) -- (storage);
\draw[arrow] (storage) -- (output);

\end{tikzpicture}
\end{center}}

\vspace{12mm}

\blueheading{Project Outcome / Deliverables (1 pts)}

\vspace{1mm}

\answerbox[3.5cm]{The project is meant to create a functional Gamified Habit Tracker. That is designed to assist users in building discipline and improving their productivity. The project will include a command line interface which enable user in creating and managing habits, the tracking of streaks, and the recording of activities. The programme will have XP rewards , bonuses for streaks , level-ups , achievements , and a health or heart mechanism. the User data will be saved locally in .dat files each time the user opens the application. The application will also provide ANSI coloured alerts and formatted terminall stats.
}

\vspace{12mm}

\blueheading{Assumptions}

\vspace{6mm}


\answerbox[2.65cm]{It is assumed that the application will have the necessary read/write file permissions to generate and reform the \verb|.dat| storage files. It is also assumed that the user is operating on terminal emulator that supports standard ANSI escape sequences for colour rendering. }
\newpage
\vspace{12mm}

\blueheading{References}

\vspace{1mm}

\answerbox[4.8cm]{(1) Bjarne Stroustrup, The C++ Programming Language. Addison-Wesley, 2013.\\
\\
(2) W. Oliveira, P. Dantas Scaico, J. Hamari, A. Soltiyeva, A. Brambilla, and S. MilHosseini, “The Effects of Gamification on Students’ Academic Performance,” Technology, Knowledge and Learning, vol. 31, no. 2, pp. 909–936, Apr. 2026, doi: 10.1007/s10758-026-09971-w.\\
\\
(3) V. W. S. Cheng, T. Davenport, D. Johnson, K. Vella, and I. B. Hickie, “Gamification in Apps and Technologies for Improving Mental Health and Well-Being: Systematic Review,” JMIR Mental Health, vol. 6, no. 6, p. e13717, June 2019, doi: 10.2196/13717.}

\end{document}

